% ==============================================================================
% Appendix E: Linear VAR Diagnostic Space and Auxiliary Specifications
% Preserves exploratory Systems 4, 5, and full-sample System 6
% ==============================================================================
\section{Linear VAR Diagnostic Space and Auxiliary Specifications}
\label{app:linear_var_diagnostics}

\subsection{Diagnostic Admissibility Criteria and Stopping Rule}
In accordance with the empirical discipline of this chapter, linear vector autoregressive models are admitted to the main text's inferential core only when they satisfy three baseline econometric criteria:
\begin{enumerate}
    \item \textbf{Modulus Stability (Criterion A):} All characteristic roots of the companion matrix lie strictly within the unit circle ($\max |\lambda_i| < 1$).
    \item \textbf{System Residual Whitening (Criterion B):} The multivariate Breusch--Godfrey Lagrange Multiplier test fails to reject the null of no residual autocorrelation at the 5\% level ($p > 0.05$).
    \item \textbf{Equation Residual Whitening (Criterion C):} Individual equation Breusch--Godfrey tests fail to reject the null of no autocorrelation at the 5\% level ($p > 0.05$) across all equations.
\end{enumerate}

During the exploratory empirical battery, three modular systems failed to satisfy these criteria: System 4 (Real Sector Dualism), System 5 (External Cost-Push), and the full-sample specification of System 6 (Central Bank Solvency). Rather than searching over ad-hoc lag specifications, non-linear trend adjustments, or arbitrary sub-samples until confirming the preferred theoretical hypotheses, the empirical audit sequence established a formal stopping rule. This appendix documents the diagnostic investigations across the lag search space ($k = 1, \dots, 12$), historical partitions (pre- and post-October 1973), and month-of-year seasonal controls, preserving the empirical record for full transparency.

\subsection{System 4: Real Sector Dualism (Descriptive Only)}
System 4 specifies the five-variable vector:
\begin{equation}
    Y^{(4)}_t = \begin{bmatrix} \pi_t & g_{H,t} & \Delta \ln \Theta_t & g_{\text{Manuf},t} & g_{\text{Mining},t} \end{bmatrix}',
\end{equation}
estimated over the full monthly sample (1960:01--1980:12, $N=250$). Table~\ref{tab:app_system4_real_dual} reports the exploratory Granger causality tests evaluated at the SBIC-selected lag ($k^*=2$).

\begin{table}[htbp]
\centering
\small
\renewcommand{\arraystretch}{0.90}
\caption{Exploratory Granger Tests: Real Sector Dualism ($N=250$, Full Sample)}
\label{tab:app_system4_real_dual}
\resizebox{\textwidth}{!}{%
\begin{tabular}{llcccc}
\toprule
\textbf{Dependent Variable} & \textbf{Predictor} & \textbf{Lag ($k^*$)} & \textbf{$F$-Statistic} & \textbf{$p$-value} & \textbf{$\sum \hat{\beta}_i$} \\
\midrule
Manufacturing Output ($g_{\text{Manuf},t}$) & Base Money Growth ($g_{H,t}$) & 2 & 3.636$^{**}$ & 0.0278 & $-0.333$ \\
Base Money Growth ($g_{H,t}$) & Manufacturing Output ($g_{\text{Manuf},t}$) & 2 & 2.306 & 0.1018 & $+0.046$ \\
Mining Output ($g_{\text{Mining},t}$) & Base Money Growth ($g_{H,t}$) & 2 & 0.908 & 0.4046 & $-0.096$ \\
Base Money Growth ($g_{H,t}$) & Mining Output ($g_{\text{Mining},t}$) & 2 & 0.654 & 0.5211 & $+0.022$ \\
Inflation ($\pi_t$) & Manufacturing Output ($g_{\text{Manuf},t}$) & 2 & 7.103$^{***}$ & 0.0010 & $-0.201$ \\
Manufacturing Output ($g_{\text{Manuf},t}$) & Structural Imbalance ($\Delta \ln \Theta_t$) & 2 & 1.526 & 0.2195 & $+0.076$ \\
\bottomrule
\end{tabular}%
}
\begin{minipage}{0.98\textwidth}
\vspace{0.3em}
\footnotesize \textbf{Notes:} Evaluated at $k^*=2$. Status: \textsc{Descriptive Only}. The system exhibits persistent residual autocorrelation (System Breusch--Godfrey $\chi^2 = 228.74, p < 0.0001$; minimum equation $p < 0.0001$). Residual autocorrelation persists when extending lags up to $k=12$, under historical calendar partitioning at October 1973, and with seasonal dummy controls. Accordingly, these statistics are not used for formal directional inference.
\end{minipage}
\end{table}

\subsection{System 5: External Cost-Push Transmission (Inconclusive)}
System 5 specifies the five-variable external transmission vector:
\begin{equation}
    Y^{(5)}_t = \begin{bmatrix} g_{P,\text{gold},t} & g_{e,t} & \Delta \ln \Theta_t & g_{H,t} & \pi_t \end{bmatrix}',
\end{equation}
estimated over 1960:01--1980:12 ($N=250$). Table~\ref{tab:app_system5_external} reports exploratory Granger tests at the SBIC-selected lag ($k^*=1$).

\begin{table}[htbp]
\centering
\small
\renewcommand{\arraystretch}{0.90}
\caption{Exploratory Granger Tests: External Cost-Push Transmission ($N=250$, Full Sample)}
\label{tab:app_system5_external}
\resizebox{\textwidth}{!}{%
\begin{tabular}{llcccc}
\toprule
\textbf{Dependent Variable} & \textbf{Predictor} & \textbf{Lag ($k^*$)} & \textbf{$F$-Statistic} & \textbf{$p$-value} & \textbf{$\sum \hat{\beta}_i$} \\
\midrule
World Gold Price ($g_{P,\text{gold},t}$) & Domestic Aggregates & 1 & \dots & $> 0.10$ & \dots \\
Inflation ($\pi_t$) & Structural Imbalance ($\Delta \ln \Theta_t$) & 1 & 5.113$^{**}$ & 0.0246 & $-0.019$ \\
Devaluation Rate ($g_{e,t}$) & Inflation ($\pi_t$) & 1 & 23.305$^{***}$ & $< 0.0001$ & $+0.689$ \\
Base Money Growth ($g_{H,t}$) & Devaluation Rate ($g_{e,t}$) & 1 & 17.037$^{***}$ & $< 0.0001$ & $+0.161$ \\
Inflation ($\pi_t$) & Devaluation Rate ($g_{e,t}$) & 1 & 1.697 & 0.1939 & $+0.046$ \\
\bottomrule
\end{tabular}%
}
\begin{minipage}{0.98\textwidth}
\vspace{0.3em}
\footnotesize \textbf{Notes:} Evaluated at $k^*=1$. Status: \textsc{Inconclusive}. The system exhibits persistent residual autocorrelation (System Breusch--Godfrey $\chi^2 = 310.16, p < 0.0001$; minimum equation $p < 0.0001$). Extending lag length to $k=12$, calendar partitioning, and seasonal controls fail to whiten the residuals. Furthermore, key relations display sensitivity to multivariate conditioning. These results are retained for documentation but not used for directional inference.
\end{minipage}
\end{table}

\subsection{System 6: Full-Sample Central Bank Solvency Diagnostics}
In the full sample (1960:01--1980:12, $N=250$), System 6 fails residual whitening across all tested lags $k \in \{1, 2, 3, 4\}$ (System Breusch--Godfrey $p < 0.0001$). However, partitioning the sample at October 1973 reveals that residual autocorrelation is concentrated in the post-1973 period under the crawling-peg regime. In the pre-October 1973 subperiod, residual autocorrelation is fully resolved at lag $k=3$ (System Breusch--Godfrey $p = 0.1262$; minimum equation $p = 0.0557$). As reported in Section~\ref{sec:granger_solvency_pre73} and Table~\ref{tab:core_granger_evidence}, System 6 is admitted to the main text strictly as a qualified, historically specific result for the pre-1973 era.
